\section*{Bab \ref{ch:b2:weather-systems} --- Sistem Cuaca}

\begin{solution}{exo:b2:weather-systems:1}
(a) $f = 2 \times 7.292 \times 10^{-5} \times \sin\ang{50} = \qty{1.117e-4}{s^{-1}}$. (b) Dengan
$\bar T = \qty{278}{K}$, $\tan\alpha = 1.117 \times 10^{-4} \times 278 \times 15/(9.81 \times 6)
= 0.00791$, sekitar 1 banding 126. (c) $1/0.00791 = 126$: sekitar \qty{126}{km} di belakang front. (d)
Kemiringan berbanding terbalik dengan kontras suhu: kemiringan itu menjadi separuhnya, 1 banding 253.
\end{solution}

\begin{solution}{exo:b2:weather-systems:2}
(a) $6 \times 150 = \qty{900}{km}$ di depan. (b) $900/25 = \qty{36}{h}$, satu setengah hari. (c) Mula-mula
sirus yang tinggi, lalu sirostratus, lalu altostratus; awan menebal dan merendah, dan hujan
turun dari awan rendah yang tebal ketika front mendekat.
\end{solution}

\begin{solution}{exo:b2:weather-systems:3}
Simbol-simbol menunjuk arah gerak front, dan ke dalam udara yang lebih hangat di depan front dingin atau
udara yang lebih dingin di depan front hangat. (a) Front dingin: udara dingin terletak di barat lautnya, di belakangnya;
front itu bergerak ke tenggara. (b) Front hangat: udara dingin terletak di utaranya, di depannya; front itu bergerak ke
utara. (c) \omterm{def:b2:weather-systems:fronts}{Front oklusi}, yang menghubungkan pusat tekanan rendah dengan titik tripel tempat front dingin dan front hangat
bermula; udara dingin terletak di kedua sisinya, dengan udara hangat terangkat di atasnya; front itu bergerak ke arah yang ditunjuk
simbol-simbolnya. (d) \omterm{def:b2:weather-systems:fronts}{Front stasioner}: segitiga-segitiganya menunjuk ke selatan, ke dalam udara hangat, sehingga udara dingin
terletak di utara; front itu nyaris tidak bergerak.
\end{solution}

\begin{solution}{exo:b2:weather-systems:4}
(a) $\beta = 2 \times 7.292 \times 10^{-5} \times \cos\ang{45}/(6.371 \times 10^6) =
\qty{1.62e-11}{m^{-1}.s^{-1}}$. (b) Untuk \qty{4000}{km}, $k = \qty{1.571e-6}{m^{-1}}$, $\beta/k^2 =
\qty{6.6}{m/s}$ dan $c = 12 - 6.6 = \qty{5.4}{m/s}$; untuk \qty{8000}{km}, $\beta/k^2 =
\qty{26.3}{m/s}$ dan $c = \qty{-14.3}{m/s}$. (c) Gelombang pendek bergerak ke timur, \qty{470}{km} sehari;
gelombang panjang bergerak ke barat, \qty{1240}{km} sehari, melawan angin.
\end{solution}

\begin{solution}{exo:b2:weather-systems:5}
(a) $c = 0$ bila $U = \beta/k^2$, sehingga $k = \sqrt{\beta/U}$ dan $\lambda = 2\pi/k =
2\pi\sqrt{U/\beta}$. (b) $2\pi\sqrt{12/(1.62 \times 10^{-11})} = \qty{5410}{km}$. (c) Lingkaran
lintang \ang{45} berukuran $2\pi \times 6371 \times \cos\ang{45} = \qty{28300}{km}$: sekitar 5.2
gelombang, jadi lima dalam praktiknya. (d) $\lambda$ tumbuh sebanding dengan $\sqrt U$: angin baratan yang lebih kuat membuat
gelombang diam lebih panjang, dan lebih sedikit gelombang yang muat mengelilingi belahan bumi.
\end{solution}

\begin{solution}{exo:b2:weather-systems:6}
(a) Pada $C_+ = (\sqrt{72}, \sqrt{72}, 27)$: $\dot x = 10(\sqrt{72} - \sqrt{72}) = 0$, $\dot y =
\sqrt{72}(28 - 27) - \sqrt{72} = 0$, $\dot z = 72 - \tfrac83 \times 27 = 72 - 72 = 0$.
(b) $\mathrm{e}^{-13.67} = 1.16 \times 10^{-6}$: volume itu menyusut hampir sejuta kali dalam satu
satuan waktu. (c) Sebuah volume dapat menyusut sementara volume itu diregangkan ke satu arah dan dimampatkan jauh lebih
kuat ke arah-arah lainnya: jarak di sepanjang arah regangan bertambah, dan dua lintasan yang berdekatan
saling menjauh di sepanjang arah itu, sementara volume di antara keduanya lenyap.
\end{solution}

\begin{solution}{exo:b2:weather-systems:7}
(a) $\ln 2/0.906 = 0.765$ satuan waktu. (b) $\ln(10/10^{-8})/0.906 = \ln(10^9)/0.906 = 22.9$ satuan
waktu. (c) $\ln 1000/0.906 = 7.6$ satuan waktu: perbaikan keadaan awal seribu kali lipat
hanya memberi sepertiga waktu yang ditemukan pada (b), karena $\ln 10^3/\ln 10^9 = 1/3$.
\end{solution}

\begin{solution}{exo:b2:weather-systems:8}
(a) $(303 - 195)/195 = 0.554$ dan $V_{\max} = \sqrt{0.9 \times 0.554 \times 15\,000} =
\qty{86.5}{m/s}$. (b) Jauh di atas \qty{70}{m/s} kategori 5, dekat dengan \qty{86}{m/s} yang
telah dilampaui lima badai belakangan ini. (c) Badai jarang bertahan lama dalam keadaan tunak yang diandaikan
teori itu; geser angin vertikal, udara kering yang tertarik ke dalam badai, dan udara yang turun di sekitarnya, yang
diabaikan oleh teori itu, melemahkan badai.
\end{solution}

\begin{solution}{exo:b2:weather-systems:9}
(a) $(303 - 195)/303 = 0.356$. (b) $(303 - 195)/195 = 0.554$; rasionya $0.554/0.356 =
303/195 = 1.554$. (c) Dalam mesin Carnot, usaha meninggalkan mesin. Dalam siklon tropis,
usaha didisipasi oleh gesekan dekat permukaan laut, dan kalor itu menghangatkan udara yang mengalir masuk pada
suhu $T_s$, tempat mesin itu menyerap kalornya: kalor itu dipakai lagi, dan mesin itu berperilaku
seolah-olah lebih efisien dengan faktor $T_s/T_o$.
\end{solution}

\begin{solution}{exo:b2:weather-systems:10}
(a) Penurunan \cref{prop:b2:weather-systems:slope} tanpa hampiran memberikan
$\mathrm{d}z/\mathrm{d}x = f(\rho_c v_c - \rho_w v_w)/[g(\rho_c - \rho_w)]$. Fluida yang lebih padat
yang bertumpu di atas fluida yang lebih ringan tidak stabil: gangguan sekecil apa pun membuatnya terbalik. Dalam front yang
diam, udara dingin yang lebih padat karena itu harus terletak di bawah, yaitu baji dalam definisi; maka
$\rho_c > \rho_w$, bidang itu naik ke arah udara dingin, dan $\rho_c v_c > \rho_w v_w$.
(b) $\zeta = \partial v/\partial x = 10/(50 \times 10^3) = \qty{2e-4}{s^{-1}}$, sekitar $1.9f$ dengan
$f = \qty{1.03e-4}{s^{-1}}$. (c) Positif di Belahan Bumi Utara: siklonik, dan sama besarnya dengan
parameter Coriolis itu sendiri; front adalah pita vortisitas siklonik yang kuat.
\end{solution}

\begin{solution}{exo:b2:weather-systems:11}
(a) $\log_2(6/0.5) = \log_2 12 = 3.58$ penggandaan, $3.58 \times 2.4 = 8.6$~hari. (b) Satu hari
prakiraan tambahan memerlukan pertumbuhan yang satu waktu penggandaan lebih pendek sebelum kesalahan menjadi tidak berguna: kesalahan
awal harus mengecil dengan faktor $2^{1/2.4} = 1.33$ setiap dasawarsa. (c) Pengurangan sepuluh kali lipat setara dengan
$\log_2 10 = 3.32$ penggandaan; jika pengurangan itu hanya memberi 4 sampai 5 hari, kesalahan kecil berlipat dua dalam $4/3.32 = 1.2$
sampai $5/3.32 = 1.5$ hari: kesalahan yang lebih kecil tumbuh lebih cepat daripada kesalahan lebih besar yang diukur dalam
prakiraan tahun 1980-an, melalui kalor laten yang dilepaskan dalam konveksi pada skala kecil.
\end{solution}

\begin{solution}{exo:b2:weather-systems:12}
(a) Pada \qty{30}{\celsius}, $e_s = \qty{43.5}{hPa}$ dan $q_s^* = \qty{27.1}{g/kg}$; udara pada
\qty{29}{\celsius} memiliki $e = 0.8 \times 41.0 = \qty{32.8}{hPa}$ dan $q_a = \qty{20.4}{g/kg}$;
$h_o^* - h^* = 1004.5 \times 1 + 2.501 \times 10^6 \times 0.0067 = \qty{17.8}{kJ/kg}$, dibandingkan dengan
\qty{16.8}{kJ/kg}. (b) $(303.15 - 200)/200 = 0.516$ dan $V_{\max} = \sqrt{0.9 \times 0.516
\times 17\,800} = \qty{90.9}{m/s}$, dibandingkan dengan \qty{87.9}{m/s}. (c) Sekitar \qty{3.4}{\percent} per
kelvin. Ketidaksetimbangan naik \qty{6}{\percent}, yang menambah \qty{3}{\percent} pada kecepatan;
faktor $(T_s - T_o)/T_o$ naik \qty{1}{\percent}, yang menambah \qty{0.5}{\percent}: bahan
bakarnya, yaitu kelembapan yang dapat diberikan laut yang lebih hangat, menyumbang paling banyak, melalui kenaikan eksponensial
$e_s$.
\end{solution}

\begin{solution}{pb:b2:weather-systems:1}
\textbf{1.} $(480 - 390)/6 = \qty{15}{km/h}$.
\textbf{2.} $390/15 = \qty{26}{h}$ setelah pukul 12 UTC: pukul 14 UTC keesokan harinya.
\textbf{3.} $f = \qty{1.117e-4}{s^{-1}}$.
\textbf{4.} Dengan $\bar T = \qty{281}{K}$, $\tan\alpha = 1.117 \times 10^{-4} \times 281 \times
16/(9.81 \times 8) = 0.0064$, sekitar 1 banding 156.
\textbf{5.} Front itu lalu berada $6 \times 15 = \qty{90}{km}$ melewati kota; \omterm{def:b2:weather-systems:fronts}{bidang front} berada
$90\,000 \times 0.0064 = \qty{576}{m}$ di atasnya.
\textbf{6.} $1500/0.0064 = \qty{234}{km}$ di belakang front, yang dicapai $234/15 = \qty{15.6}{h}$ setelah
front lewat.
\textbf{7.} Kumulus dan kumulonimbus, hujan deras sesaat, mungkin guntur, embusan angin, dan penurunan
suhu yang tiba-tiba; tekanan, yang turun ketika front mendekat, mulai naik ketika front lewat; lalu
udara yang lebih sejuk dan lebih cerah.
\textbf{8.} $\beta = 2 \times 7.292 \times 10^{-5} \times \cos\ang{50}/(6.371 \times 10^6) =
\qty{1.47e-11}{m^{-1}.s^{-1}}$.
\textbf{9.} $k = 2\pi/(5 \times 10^6) = \qty{1.257e-6}{m^{-1}}$ dan $\beta/k^2 = \qty{9.3}{m/s}$.
\textbf{10.} $c = 20 - 9.3 = \qty{10.7}{m/s}$, ke arah timur.
\textbf{11.} $10.7 \times 86\,400 = \qty{924}{km}$ ke timur.
\textbf{12.} $\lambda_s = 2\pi\sqrt{20/(1.47 \times 10^{-11})} = \qty{7330}{km}$; lingkaran
lintang \ang{50} berukuran $2\pi \times 6371 \times \cos\ang{50} = \qty{25700}{km}$, sekitar 3.5
gelombang seperti itu.
\textbf{13.} Di depannya, di sisi timurnya, tempat gelombang pendek memicu udara naik dan
tekanan permukaan turun.
\textbf{14.} $\lambda = \ln(3.9 \times 10^{-3}/5 \times 10^{-7})/(14 - 4) = \ln 7800/10 = 0.896$.
\textbf{15.} $\ln 2/0.896 = 0.77$ satuan waktu.
\textbf{16.} \qty{1}{\percent} di bawah 0.906: rata-rata 200 pasang, waktu yang terbatas, dan pembacaan dari
grafik.
\textbf{17.} Begitu kedua simulasi berjauhan, keduanya terletak pada titik-titik atraktor yang tidak berkaitan, dan
jaraknya tidak dapat melebihi ukuran atraktor: pertumbuhannya jenuh.
\textbf{18.} $\ln 10/0.896 = 2.6$ satuan waktu kemudian.
\textbf{19.} $\log_2(3/0.1) = \log_2 30 = 4.91$ penggandaan.
\textbf{20.} $4.91 \times 2.4 = 11.8$~hari.
\textbf{21.} Satu waktu penggandaan: 2.4~hari.
\textbf{22.} $4.91 \times 5 = 24.5$~hari: perkiraan awal itu jauh terlalu optimistis.
\textbf{23.} Dengan waktu penggandaan 2.4 hari, kesalahan yang sepuluh kali lebih kecil akan memberi tambahan $3.32 \times 2.4
= 8.0$~hari, bukan 4 sampai 5: kesalahan kecil tumbuh lebih cepat daripada 2.4 hari per penggandaan, dan
perkiraan yang dibangun di atas waktu penggandaan itu melebih-lebihkan cakrawalanya.
\textbf{24.} Prakiraan ini dapat berguna paling lama sekitar \textbf{12 hari}, dan dalam praktiknya
kurang dari itu.
\end{solution}
